% GENERATED FILE. Edit canonical lessons, then run npm run generate:books.
% canonical-source-hash: fnv1a64:b8b921879393aee7
% canonical-lessons: FA-C199-alude, FA-C199-charb, FA-C199-qavi, FA-C199-rahat, FA-C199-shik

\chapter{Polluted, Greasy, Strong, Comfortable, Smart}
\label{ch:fa-polluted-greasy-strong-comfortable-smart}
\hlchaptermodality{199}

\begin{chapteropening}
\textbf{By the end of this chapter:} \emph{I can say polluted, greasy, strong, comfortable and smart.}

Complete the last lesson of chapter 199: I can say polluted, greasy, strong, comfortable and smart.
\end{chapteropening}

\section[\texorpdfstring{\fa{آلوده} (âlude)}{آلوده (âlude)}]{\fa{آلوده} (âlude) — polluted, contaminated}

\label{lesson:FA-C199-alude}



\begin{quote}
Before the new one: say the Persian for an example, then the Persian for success.
\end{quote}

\subsection*{You'll want to know: \fa{آلوده}}
\textbf{\fa{آلوده}} — \emph{âlude} — \textquotedblleft{}polluted, contaminated\textquotedblright{}.

\begin{cousinweb}
\emph{âlude} is the past participle of \emph{âludan}, \textquotedblleft{}to soil, to stain\textquotedblright{}: \textquotedblleft{}stained\textquotedblright{}. The verb is native, usually traced to Middle Persian \emph{ālūdan}.
\end{cousinweb}

\begin{grammarlens}[title={What you've built}]
One more word for this chapter.
\end{grammarlens}

\subsection*{Your turn}
\begin{itemize}
\raggedright
  \item \emph{Say it:} \emph{âlude}
  \item \emph{Say it:} \emph{âlude}, once more
  \item \emph{Say it:} say \emph{âlude}
  \item \emph{Recall:} say the Persian for an example, then the Persian for success, then say \emph{âlude} again
\end{itemize}

\begin{tcolorbox}[breakable,colback=teal!4,colframe=teal!35!black,title={Before you move on}]
What does \fa{آلوده} mean? (\textquotedblleft{}Polluted, contaminated\textquotedblright{}.) Say it once more.
\end{tcolorbox}
\section[\texorpdfstring{\fa{چرب} (charb)}{چرب (charb)}]{\fa{چرب} (charb) — greasy, fatty}

\label{lesson:FA-C199-charb}



\begin{quote}
Before the new one: say the Persian for success, then the Persian for polluted.
\end{quote}

\subsection*{You'll want to know: \fa{چرب}}
\textbf{\fa{چرب}} — \emph{charb} — \textquotedblleft{}greasy, fatty\textquotedblright{}.

\begin{cousinweb}
\emph{charb} is a native Persian word, inherited rather than borrowed: it continues Middle Persian \emph{čarb}, \textquotedblleft{}fatty\textquotedblright{} (EtymDB).
\end{cousinweb}

\begin{grammarlens}[title={What you've built}]
One more word for this chapter.
\end{grammarlens}

\subsection*{Your turn}
\begin{itemize}
\raggedright
  \item \emph{Say it:} \emph{charb}
  \item \emph{Say it:} \emph{charb}, once more
  \item \emph{Say it:} say \emph{charb}
  \item \emph{Recall:} say the Persian for success, then the Persian for polluted, then say \emph{charb} again
\end{itemize}

\begin{tcolorbox}[breakable,colback=teal!4,colframe=teal!35!black,title={Before you move on}]
What does \fa{چرب} mean? (\textquotedblleft{}Greasy, fatty\textquotedblright{}.) Say it once more.
\end{tcolorbox}
\section[\texorpdfstring{\fa{قوی} (qavi)}{قوی (qavi)}]{\fa{قوی} (qavi) — strong}

\label{lesson:FA-C199-qavi}



\begin{quote}
Before the new one: say the Persian for polluted, then the Persian for greasy.
\end{quote}

\subsection*{You'll want to know: \fa{قوی}}
\textbf{\fa{قوی}} — \emph{qavi} — \textquotedblleft{}strong\textquotedblright{}.

\begin{cousinweb}
\emph{qavi} is borrowed from Arabic \emph{qawī}, \textquotedblleft{}strong\textquotedblright{}. Its Arabic root is \emph{q-w-y} (Lane).
\end{cousinweb}

\begin{grammarlens}[title={What you've built}]
One more word for this chapter.
\end{grammarlens}

\subsection*{Your turn}
\begin{itemize}
\raggedright
  \item \emph{Say it:} \emph{qavi}
  \item \emph{Say it:} \emph{qavi}, once more
  \item \emph{Say it:} say \emph{qavi}
  \item \emph{Recall:} say the Persian for polluted, then the Persian for greasy, then say \emph{qavi} again
\end{itemize}

\begin{tcolorbox}[breakable,colback=teal!4,colframe=teal!35!black,title={Before you move on}]
What does \fa{قوی} mean? (\textquotedblleft{}Strong\textquotedblright{}.) Say it once more.
\end{tcolorbox}
\section[\texorpdfstring{\fa{راحت} (râhat)}{راحت (râhat)}]{\fa{راحت} (râhat) — comfortable}

\label{lesson:FA-C199-rahat}



\begin{quote}
Before the new one: say the Persian for greasy, then the Persian for strong.
\end{quote}

\subsection*{You'll want to know: \fa{راحت}}
\textbf{\fa{راحت}} — \emph{râhat} — \textquotedblleft{}comfortable\textquotedblright{}.

\begin{cousinweb}
\emph{râhat} is borrowed from Arabic \emph{rāḥa}, \textquotedblleft{}rest, ease\textquotedblright{}. Its Arabic root is \emph{r-w-ḥ} (Lane).
\end{cousinweb}

\begin{grammarlens}[title={What you've built}]
One more word for this chapter.
\end{grammarlens}

\subsection*{Your turn}
\begin{itemize}
\raggedright
  \item \emph{Say it:} \emph{râhat}
  \item \emph{Say it:} \emph{râhat}, once more
  \item \emph{Say it:} say \emph{râhat}
  \item \emph{Recall:} say the Persian for greasy, then the Persian for strong, then say \emph{râhat} again
\end{itemize}

\begin{tcolorbox}[breakable,colback=teal!4,colframe=teal!35!black,title={Before you move on}]
What does \fa{راحت} mean? (\textquotedblleft{}Comfortable\textquotedblright{}.) Say it once more.
\end{tcolorbox}
\section[\texorpdfstring{\fa{شیک} (shik)}{شیک (shik)}]{\fa{شیک} (shik) — smart, stylish}

\label{lesson:FA-C199-shik}



\begin{quote}
Before the new one: say the Persian for strong, then the Persian for comfortable.
\end{quote}

\subsection*{You'll want to know: \fa{شیک}}
\textbf{\fa{شیک}} — \emph{shik} — \textquotedblleft{}smart, stylish\textquotedblright{}.

\begin{cousinweb}
\emph{shik} is borrowed from French \emph{chic}, \textquotedblleft{}elegant\textquotedblright{} (EtymDB).
\end{cousinweb}

\begin{grammarlens}[title={What you've built}]
One more word for this chapter.
\end{grammarlens}

\subsection*{Your turn}
\begin{itemize}
\raggedright
  \item \emph{Say it:} \emph{shik}
  \item \emph{Say it:} \emph{shik}, once more
  \item \emph{Say it:} say \emph{shik}
  \item \emph{Recall:} say the Persian for strong, then the Persian for comfortable, then say \emph{shik} again
\end{itemize}

\begin{tcolorbox}[breakable,colback=teal!4,colframe=teal!35!black,title={Before you move on}]
What does \fa{شیک} mean? (\textquotedblleft{}Smart, stylish\textquotedblright{}.) Say it once more.
\end{tcolorbox}
